% =====================================================================
% 5. Huellas de edificios (semana 3)
% =====================================================================
\section{Huellas de edificios}
\label{sec:edificios}

\campo{make edificios} (\campo{src/edificios/run.py}) procesa cada fuente
con el mismo esquema: descarga con SHA-256, prefiltro a PDET por la regla del
centroide, limpieza compartida (\campo{comun.py}) hacia un Parquet con el
mismo contrato de columnas, y carga e índices en MongoDB.

\begin{table}[H]
  \centering\small
  \caption{Adquisición y limpieza de las dos fuentes.}
  \label{tab:edif-resumen}
  \begin{tabularx}{\textwidth}{@{}L r r@{}}
    \toprule
    & \textbf{Google Open Buildings v3} & \textbf{Microsoft Building Footprints} \\
    \midrule
    Archivos descargados & 8 celdas S2 (3,11\,GB) & 125 \emph{quadkeys} (373,4\,MB) \\
    Edificios leídos & 16\,512\,141 (en Colombia) & 4\,595\,615 \\
    Edificios en PDET (en la colección) & 2\,694\,786 & 1\,761\,048 \\
    Polígonos inválidos reparados (GEOS) & 0 & 28 \\
    Reparados en la carga (validación esférica) & 0 & 29 \\
    Duplicados por \campo{geom\_hash} & 0 & 0 \\
    Licencia & CC BY 4.0 / ODbL & ODbL \\
    \bottomrule
  \end{tabularx}
\end{table}

\subsection{Google Open Buildings: descarga y prefiltro}
\label{sec:google}

Google publica sus edificios en CSV comprimidos, uno por celda S2 de nivel~4.
Cada edificio trae su polígono en WKT, el centroide, el área que calcula
Google (\campo{area\_in\_meters}) y la confianza del modelo
(\campo{confidence}, 0,65 a 1,0).

\textbf{Descarga.} \campo{descarga\_google.py} cruza la tabla de celdas que
publica Google con la unión de los 170 municipios: quedan 8 celdas (3,11\,GB)
de las 10 que tocan Colombia. Cada celda se baja en \emph{streaming}, se
reanuda si se corta, se verifica su tamaño y se calcula su SHA-256; el SHA-256
de la fuente (hash de la lista ordenada de pares \campo{token:sha256}) se
guarda en \campo{fuentes} y dio el mismo valor en dos máquinas distintas.

\textbf{Prefiltro con DuckDB.} Las 8 celdas tienen 31 millones de filas en el
rectángulo de Colombia, así que no se cargan completas: DuckDB lee cada CSV
comprimido en tandas de unas 100\,000 filas y descarta en SQL lo que cae fuera
de los rectángulos PDET. Luego se prueba con un índice espacial de Shapely si
el centroide cae dentro de algún municipio PDET. De 31\,370\,317 filas en el
rectángulo, 16\,512\,141 están dentro de Colombia y 2\,694\,786 en los
municipios PDET (16,3\,\%). Una corrida de control con las 10 celdas dio
exactamente los mismos 2\,694\,786 edificios.

\fig{cap_google_prefiltro.png}{0.28}{Prefiltro de Google: acumulados en
Colombia y tras el prefiltro por cada celda.}{fig:g-prefiltro}

\subsection{Microsoft Building Footprints: descarga y prefiltro}
\label{sec:microsoft}

Microsoft publica sus huellas globales (versión 2026-02-03) como GeoJSON por
líneas comprimido, en particiones por \emph{quadkey} de nivel 9; solo traen el
polígono. \campo{descarga\_microsoft.py} lee el índice público
\campo{dataset-links.csv}, toma las 232 particiones de Colombia, se queda con
las 125 que intersectan la unión de los municipios PDET y las baja en paralelo
(8 hilos), con un \campo{manifest.csv} que guarda el SHA-256 de cada archivo y
uno global. \campo{limpieza\_microsoft.py} procesa cada archivo en un proceso
aparte (8 en paralelo) y conserva solo los edificios cuyo centroide cae en
PDET: de 4\,595\,615 quedan 1\,761\,048.

\fig{cap_ms_limpieza.png}{0.25}{Tabla de control de la limpieza de
Microsoft.}{fig:ms-limpieza}

\subsection{Limpieza común y contrato del Parquet}
\label{sec:limpieza}

Las dos fuentes pasan por la misma limpieza, para que los dos Parquet
intermedios tengan las mismas columnas:

\begin{enumerate}[leftmargin=*]
  \item Se construyen las geometrías con Shapely y se descartan las ilegibles.
  \item \textbf{Validación y reparación:} los polígonos inválidos se reparan
        con \campo{make\_valid}; como el contrato pide \campo{Polygon}, se
        conserva la parte poligonal mayor.
  \item \textbf{Área y centroide:} se proyecta a EPSG:9377 para calcular
        \campo{area\_m2} y el centroide; la geometría se guarda en EPSG:4326.
  \item \textbf{\campo{geom\_hash}:} hash de la geometría normalizada y
        redondeada, para que dos copias del mismo edificio den el mismo valor;
        se descartan los repetidos (un edificio puede venir en dos archivos).
  \item DuckDB une las partes en un único Parquet comprimido (zstd) con
        metadatos GeoParquet.
\end{enumerate}

\begin{table}[H]
\centering\small
\caption{Contrato de columnas de los Parquet intermedios.}
\label{tab:contrato}
\begin{tabularx}{\textwidth}{@{}p{3.6cm}p{2.2cm}L@{}}
\toprule
\textbf{Columna} & \textbf{Tipo} & \textbf{Contenido} \\
\midrule
\campo{fuente\_id} & texto & \campo{google} o \campo{microsoft} \\
\campo{geometry} & WKB & \campo{Polygon} en EPSG:4326 \\
\campo{centroide\_lon}, \campo{centroide\_lat} & double & Centroide de la huella, en EPSG:4326 \\
\campo{area\_m2} & double & Área en EPSG:9377, en m\textsuperscript{2} \\
\campo{area\_m2\_fuente}, \campo{confianza} & double o null & Área y confianza de Google; null en Microsoft \\
\campo{geom\_hash} & texto & Hash de la geometría normalizada \\
\bottomrule
\end{tabularx}
\end{table}

\subsection{Carga eficiente en MongoDB}
\label{sec:carga}

Con 4,5 millones de documentos, la eficiencia de la carga depende de no
cargar todo a memoria y de no mantener índices mientras se inserta.
\campo{carga.py}:

\begin{itemize}
  \item Lee el Parquet con \campo{pyarrow} en \textbf{lotes de 5\,000 filas}
        (\campo{config.yaml}); la geometría pasa de WKB a GeoJSON con Shapely.
  \item Inserta con \campo{insert\_many(ordered=False)}: un documento
        rechazado no frena el resto del lote.
  \item \textbf{Crea los índices después de cargar}, porque mantenerlos
        durante la inserción es más lento. La excepción es el índice de
        \campo{geometry} (sección~\ref{sec:s2}).
  \item Es \textbf{idempotente}: borra los documentos previos de la fuente
        antes de cargar; correrla dos veces deja el mismo conteo.
  \item Marca cada edificio con un \campo{lote\_id}, que es el \campo{\_id} de
        su fila en \campo{bitacora\_cargas}, con inicio, fin, registros y
        velocidad.
\end{itemize}

\begin{table}[H]
\centering\small
\caption{Carga a MongoDB (MongoDB 8 en Docker local, un solo equipo; los
tiempos dependen de la máquina).}
\label{tab:carga}
\begin{tabularx}{\textwidth}{@{}lrrrrL@{}}
\toprule
\textbf{Fuente} & \textbf{Documentos} & \textbf{Rechazados} & \textbf{Reparados (S2)} & \textbf{Tiempo (s)} & \textbf{Velocidad} \\
\midrule
Google    & 2\,694\,786 & 0 & 0  & 56,6 & 47\,621 docs/s \\
Microsoft & 1\,761\,048 & 0 & 29 & 37,7 & 46\,753 docs/s \\
\bottomrule
\end{tabularx}
\end{table}

\begin{table}[H]
\centering\small
\caption{Tiempo de creación de los índices de edificios.}
\label{tab:indices-edif}
\begin{tabularx}{\textwidth}{@{}p{3.6cm}Lrr@{}}
\toprule
\textbf{Índice} & \textbf{Para qué} & \textbf{Google (s)} & \textbf{Microsoft (s)} \\
\midrule
\campo{2dsphere} en \campo{centroide} & \campo{\$geoWithin} de la asignación & 6,9 & 4,4 \\
\campo{2dsphere} en \campo{geometry} & \campo{\$geoIntersects} en límites & \multicolumn{2}{c}{creado antes de cargar} \\
\campo{cod\_mpio} & agregación por municipio & 2,2 & 1,4 \\
\campo{geom\_hash} (único) & impedir duplicados & 3,4 & 2,3 \\
\bottomrule
\end{tabularx}
\end{table}

En disco, \campo{edificios\_google} ocupa 1\,279\,MB de datos y 245\,MB de
índices, y \campo{edificios\_microsoft} 835\,MB y 175\,MB. El índice único
sobre \campo{geom\_hash} es una segunda barrera contra duplicados: si alguno
llegara a la colección, la creación del índice fallaría.

\subsection{Un hallazgo: polígonos válidos para GEOS que MongoDB rechaza}
\label{sec:s2}

Con todos los índices creados al final, el \campo{2dsphere} sobre
\campo{geometry} de Microsoft falló: MongoDB construye sus claves con geometría
esférica (biblioteca S2) y encontró un polígono con aristas que se cruzan,
aunque GEOS, que trabaja en el plano, lo había dado por válido. Con una sola
geometría así, la construcción de todo el índice aborta. En una muestra de
300\,000 edificios hubo 5 casos, unos 2 por millón.

La solución fue crear ese índice \emph{antes} de cargar: así MongoDB rechaza
solo los documentos afectados y la carga sigue. Cada rechazado se repara
redondeando las coordenadas a $10^{-6}$ grados (unos 11\,cm) con
\campo{shapely.set\_precision} y se reintenta; \campo{buffer(0)} y
\campo{simplify} no resolvieron ninguno de los casos de prueba. En la carga
completa, Microsoft tuvo 29 reparados y Google ninguno, y no se descartó ningún
edificio.

\subsection{Evidencia de que los índices se usan}
\label{sec:evidencia-edif}

\campo{evidencia.py} corre la consulta de asignación sobre el polígono de
Caldono (19137) con el índice y forzando el recorrido completo
(\campo{hint(\{\$natural: 1\})}).

\begin{table}[H]
\centering\small
\caption{Planes de ejecución en Caldono (una consulta por caso; tiempos
indicativos).}
\label{tab:evidencia-edif}
\begin{tabularx}{\textwidth}{@{}llLrr@{}}
\toprule
\textbf{Fuente} & \textbf{Plan} & \textbf{Consulta} & \textbf{Resultados} & \textbf{Tiempo (s)} \\
\midrule
Google    & IXSCAN   & \campo{\$geoWithin} sobre \campo{centroide}   & 15\,303 & 0,44 \\
Google    & IXSCAN   & \campo{\$geoIntersects} sobre \campo{geometry} & 15\,304 & 0,37 \\
Google    & COLLSCAN & \campo{\$geoWithin} sin índice                 & 15\,303 & 2,93 \\
Microsoft & IXSCAN   & \campo{\$geoWithin} sobre \campo{centroide}   & 12\,547 & 0,32 \\
Microsoft & IXSCAN   & \campo{\$geoIntersects} sobre \campo{geometry} & 12\,548 & 0,28 \\
Microsoft & COLLSCAN & \campo{\$geoWithin} sin índice                 & 12\,547 & 1,93 \\
\bottomrule
\end{tabularx}
\end{table}

La diferencia de un edificio entre \campo{\$geoWithin} (centroide dentro) y
\campo{\$geoIntersects} (huella que toca el municipio) es un edificio que cruza
el límite: es la razón de usar el centroide, para que cada edificio pertenezca
a un solo municipio.
